\documentclass[12pt]{article}

\usepackage{fontawesome}
\usepackage{hyperref}
\usepackage{xurl}
\usepackage[tagged, highstructure]{accessibility}
\usepackage{bookmark}
\usepackage{enumitem}

\hypersetup{
    colorlinks=false,
    pdfborder={0 0 0},
    pdftitle={More Commands},
    pdfauthor={Adrianna Holden-Gouveia},
    pdfsubject={Introduction to Linux - Lab Assignment},
    pdflang={en-US},
    pdfkeywords={lab assignment, More Commands}
}

% Configure PDF bookmarks for navigation
\bookmarksetup{
    numbered,
    open,
}

% Configure list spacing for better accessibility
\setlist{nosep}



\title{More Commands}
\author{
        Adrianna Holden-Gouveia \\
        Website: \url{https://aholdengouveia.name}\\ 
        \date{\vspace{-5ex}}
        %Email: \href{mailto:admin@aholdengouveia.name}{admin@aholdengouveia.name} \\
%        \faLinkedin{: aholdengouveia} \\
        \faGithub {: aholdengouveia} \\
%        \faTwitter {: aholdengouveia} \\
        }

%S\date{\today}


\begin{document}    

\maketitle

\section*{Accessibility Notice}
This document is also available in HTML format at:

\url{https://aholdengouveia.name/IntroLinux/labs/cmdsandwildcards.html}

The HTML version provides enhanced accessibility features including keyboard navigation, screen reader support, responsive design, dark mode support, and high contrast options.


%\begin{abstract}
%This is a template for Linux Administration Lab work
%\end{abstract}
%\tableofcontents

\section*{Objectives:}
\begin{itemize}
    \item Learn to use wildcards for file pattern matching
    \item Practice using the ls command with various wildcard patterns
    \item Understand the difference between ls and find commands
    \item Develop skills in efficiently locating files in Linux
\end{itemize}
\section*{Complete the following problems}

References, a video, a PowerPoint and some notes are available at my website
\url {https://www.aholdengouveia.name/IntroLinux/cmdsandwildcards.html}

\subsection*{How to setup for this lab}

\begin{itemize}
    \item We are going to create a test directory with numerous files and then use find and wildcards to locate files.

    \item Create a test subdirectory for this lab. I do NOT recommend using any spaces in the name of the files or folders.
    \item Using touch create at least a dozen files with different names with the following characteristics:
        \begin{enumerate}
            \item Some start with f
            \item Some start with F
            \item Some start with a number
            \item Some end with a number
            \item Some do not have an f or F
            \item Some do not have a number
            \item At least one starts with a . (hidden file)
        \end{enumerate}

    \item If you're struggling to think of names, use a naming scheme such as characters from your favorite video game, or locations and characters from favorite books etc.  Make sure your files are full names of at least 5 characters, not just single letters/numbers
    \item For each of the problems you should start with ls and then figure out the pattern needed to show the requested information, then pick 1 of them to try with find instead.  \item Demo of how to use ls is in the Linux FAQ videos \url{https://youtube.com/playlist?list=PLQOmuLC45eoQlIMH7oCNL6yFP7kWQubBo&si=R1SFWzEPHDDo-VBJ}

    \end{itemize}
\subsection*{Problems to solve}
Use wildcards and the ls command to answer the following. 
    \begin{enumerate}
        \item All files
        \item All files including hidden files.
        \item Files that start with “f”
        \item Files that start with “f” or “F”
        \item Files that have “f” or “F” anywhere in the name
        \item Files that have a digit in the name
        \item Files that start with a through d
        \item Files that end with 1 or 2
        \item Experiment with the find command by solving one of the above questions using find instead of ls.

    \end{enumerate}



\section*{Deliverables}
One file that has the lab title, name and date and the answers.  
\begin{itemize}
    \item Each problem given should be numbered to match the question
    \item Make sure to list each command and the result.  \item Make sure to take a screenshot of each successful command.  
    \item Don't use one screenshot for everything, make sure each screenshot is labeled with which problem it solved.
\end{itemize}
\end{document}